\begin{statementbox}
\textbf{定义。} 过焦点且垂直于椭圆长轴、双曲线实轴或抛物线对称轴的\textbf{弦段}称为通径。
记整条通径长为 $L$，则
\[
 \boxed{L_{\text{椭、双}}=\frac{2b^2}{a}},\qquad
 \boxed{L_{\text{抛}}=2p}.
\]
椭圆 $a>b>0$，$a,b$ 为长、短半轴；双曲线 $a,b>0$，分别为实、虚半轴。
抛物线写成 $y^2=\pm2px$ 或 $x^2=\pm2py$，$p>0$ 是焦点到准线的距离。
\medskip
\textbf{更推荐的速记写法（系数 $k$）：}
\[
 \boxed{y^2=kx\ (k\ne0):\quad
 F=\left(\frac{k}{4},0\right),\quad
 \text{准线 }x=-\frac{k}{4},\quad
 L=|k|.}
\]
当 $k>0$（向右开口）时，直接记作
\[
 \boxed{\text{焦点横坐标}=\frac{k}{4},\quad
 \text{准线横坐标}=-\frac{k}{4},\quad
 \text{整通径}=k.}
\]
它与原写法并不矛盾：$p=|k|/2$，因此 $L=2p=|k|$；
\textbf{半通径}为 $|k|/2$，不是 $|k|$。
若写成 $x^2=ky$，则 $F=(0,k/4)$、准线 $y=-k/4$、$L=|k|$。
横向有心曲线的焦点为 $(\pm c,0)$，相应通径端点是
\[
 (c,\pm b^2/a)\quad\text{或}\quad(-c,\pm b^2/a).
\]
右开口抛物线端点为 $(p/2,\pm p)$，左开口为 $(-p/2,\pm p)$；纵向交换坐标。
$b^2/a$、$p$ 是\textbf{半通径长}，整条弦须乘 $2$。
平移或旋转不改变该长度关系，但端点坐标要随坐标位置调整。


\medskip
\textbf{考场识别：}先确认“经过焦点”与“垂直焦点所在的轴”两个条件，再套通径式；
若仅经过焦点但与轴不垂直，那是一般焦点弦，不能直接套本公式。

\textbf{把公式写成“半长 $\times2$”：}
\[
 \underbrace{\frac{L}{2}}_{\text{焦点到端点的距离}}
 =\begin{cases}
 \dfrac{b^2}{a},&\text{椭圆、双曲线},\\[2pt]
 p,&\text{抛物线}.
 \end{cases}
\]
其中焦点恰是通径的中点（依赖曲线关于焦点所在轴的对称性）。
对于纵向椭圆、纵向双曲线，交换上述端点坐标的两个分量即可。
\textbf{注意：}焦点在曲线的对称轴上，但“通径本身”与该轴垂直，二者不是同一条直线。

\end{statementbox}
